\chapter{Convex and Concave Functions}

A convex set has no dents: it contains the whole chord between any two of
its points. This chapter transfers that idea from sets to functions. A
function is \emph{concave} if every chord of its graph lies on or below the
graph, \emph{convex} if every chord lies on or above it, the graph bends
consistently one way. The technical engine, the Hessian and its
definiteness, was drilled in the companion volumes (\emph{Derivatives \&
Partial Derivatives} and \emph{Matrix Algebra}); what is new here is
\emph{where} we evaluate it. Definiteness at a single critical point
classifies that point locally; semidefiniteness at \emph{every} point of the
domain classifies the whole function globally.

The reward is the central convenience of optimization theory: for a concave
objective the first-order conditions are not merely necessary but
\emph{sufficient}, and under strict concavity the maximizer is unique. We
close with the weaker ordinal notion of quasiconcavity (convex upper
contour sets) which survives increasing transformations and is the
natural home of the Cobb--Douglas function: concave only when its exponents
sum to at most one, but quasiconcave always.

\section{Core Rules}

\begin{rulebox}{The chord criterion}
Throughout, the domain $S\subseteq\R^{n}$ is a \emph{convex set}:
otherwise a chord could leave the domain. A function $f:S\to\R$ is
\begin{itemize}[leftmargin=1.6em,itemsep=0.3ex]
  \item \emph{convex} if for all $\mathbf{x},\mathbf{y}\in S$
        and all $\lambda\in[0,1]$
        \[
          f\bigl(\lambda\mathbf{x}+(1-\lambda)\mathbf{y}\bigr)
          \;\le\; \lambda f(\mathbf{x})+(1-\lambda)f(\mathbf{y})
        \]
        (the chord lies on or above the graph);
  \item \emph{concave} if the reverse inequality $\ge$ holds
        (the chord lies on or below the graph).
\end{itemize}
\emph{Strictly} convex/concave: the inequality is strict whenever
$\mathbf{x}\neq\mathbf{y}$ and $\lambda\in(0,1)$. Basic facts: $f$ is
concave iff $-f$ is convex; affine functions
$f(\mathbf{x})=\mathbf{a}^{T}\mathbf{x}+b$ are both convex and concave, and
neither strictly.
\end{rulebox}

\begin{rulebox}{One variable: the second-derivative test}
Let $f$ be twice differentiable on an interval $I$. Then
\[
  f \text{ convex on } I \iff f''(x)\ge0 \text{ on } I,
  \qquad
  f \text{ concave on } I \iff f''(x)\le0 \text{ on } I .
\]
Moreover $f''>0$ on $I$ $\Rightarrow$ strictly convex, and $f''<0$
$\Rightarrow$ strictly concave: sufficient but \emph{not} necessary:
$f(x)=x^{4}$ is strictly convex although $f''(0)=0$.
\end{rulebox}

\begin{rulebox}{Several variables: the Hessian test}
Let $S\subseteq\R^{n}$ be open and convex and $f\in C^{2}(S)$. Then
\begin{itemize}[leftmargin=1.6em,itemsep=0.3ex]
  \item $f$ is convex on $S$ $\iff$ $\Hess_f(\mathbf{x})$ is positive
        \emph{semi}definite (PSD) for every $\mathbf{x}\in S$;
  \item $f$ is concave on $S$ $\iff$ $\Hess_f(\mathbf{x})$ is negative
        \emph{semi}definite (NSD) for every $\mathbf{x}\in S$;
  \item $\Hess_f$ PD everywhere $\Rightarrow$ $f$ strictly convex;
        $\Hess_f$ ND everywhere $\Rightarrow$ $f$ strictly concave.
        (Sufficient, not necessary: $-x^{4}-y^{4}$ is strictly concave with
        a singular Hessian at the origin.)
\end{itemize}
The contrast with the local theory: definiteness \emph{at one critical
point} classifies that point; semidefiniteness \emph{at every point of $S$}
classifies the function.
\end{rulebox}

\begin{rulebox}{Checking semidefiniteness: the minor trap}
Sylvester's leading-minor criterion decides the \emph{strict} cases: for
$H=\begin{pmatrix}a&b\\b&c\end{pmatrix}$, $D_1=a$, $D_2=ac-b^{2}$,
\[
  D_1>0,\ D_2>0 \Rightarrow \text{PD}; \qquad
  D_1<0,\ D_2>0 \Rightarrow \text{ND}; \qquad
  D_2<0 \Rightarrow \text{indefinite}.
\]
For \emph{semi}definiteness the leading minors are not enough
($\operatorname{diag}(0,-1)$ has $D_1=D_2=0$ but is NSD, not PSD); one must
check \emph{all} principal minors. Practical $2\times2$ rule:
\[
  H \text{ PSD} \iff a\ge0,\ c\ge0,\ ac-b^{2}\ge0;
  \qquad
  H \text{ NSD} \iff a\le0,\ c\le0,\ ac-b^{2}\ge0 .
\]
In general: PSD $\iff$ every principal minor is $\ge0$; NSD $\iff$ every
principal minor of order $k$ has sign $(-1)^{k}$ or vanishes.
\end{rulebox}

\begin{rulebox}{Operations that preserve concavity}
Let $f,g$ be concave on a convex $S$.
\begin{itemize}[leftmargin=1.6em,itemsep=0.3ex]
  \item \emph{Nonnegative weighted sums:} $\alpha f+\beta g$ is concave for
        $\alpha,\beta\ge0$; strictly concave if in addition one summand is
        strictly concave and carries a positive weight.
  \item \emph{Pointwise minimum:} $\min\{f,g\}$ (indeed the pointwise
        infimum of any family of concave functions) is concave. Mirror
        image: the pointwise maximum of convex functions is convex.
  \item \emph{Concave increasing composition:} if $g$ is concave and
        $F:\R\to\R$ is increasing \emph{and concave}, then $F\circ g$ is
        concave. (Convex version: $F$ increasing and convex, $g$ convex
        $\Rightarrow$ $F\circ g$ convex.)
  \item \emph{Affine substitution:} if $f$ is concave, so is
        $\mathbf{x}\mapsto f(A\mathbf{x}+\mathbf{b})$.
\end{itemize}
Warning: products of concave functions and general increasing transforms do
\emph{not} preserve concavity ($\mathrm{e}^{-x^{2}-y^{2}}$ is an increasing
transform of a concave function but is not concave, Problem~17). What an
increasing transform does preserve is quasiconcavity.
\end{rulebox}

\begin{rulebox}{Quasiconcavity and quasiconvexity}
For $c\in\R$ the \emph{upper contour set} and
\emph{lower contour set} of $f:S\to\R$ are
\[
  U_c=\{\mathbf{x}\in S: f(\mathbf{x})\ge c\}, \qquad
  L_c=\{\mathbf{x}\in S: f(\mathbf{x})\le c\}.
\]
$f$ is \emph{quasiconcave} iff every $U_c$ is convex,
equivalently
\[
  f\bigl(\lambda\mathbf{x}+(1-\lambda)\mathbf{y}\bigr)
  \;\ge\; \min\{f(\mathbf{x}),f(\mathbf{y})\} ;
\]
\emph{strictly} quasiconcave if the inequality is strict for
$\mathbf{x}\neq\mathbf{y}$, $\lambda\in(0,1)$. $f$ is \emph{quasiconvex}
iff every $L_c$ is convex (equivalently $f\le\max\{\cdot,\cdot\}$ along
chords), iff $-f$ is quasiconcave. Facts:
\begin{itemize}[leftmargin=1.6em,itemsep=0.3ex]
  \item concave $\Rightarrow$ quasiconcave (convex $\Rightarrow$
        quasiconvex); the converse fails.
  \item if $F$ is increasing and $f$ quasiconcave, then $F\circ f$ is
        quasiconcave: $F$ merely relabels the level curves, leaving every
        contour set unchanged. Quasiconcavity is an \emph{ordinal}
        property.
  \item every monotone function of one variable is both quasiconcave and
        quasiconvex (all contour sets are intervals).
\end{itemize}
\end{rulebox}

\begin{rulebox}{Cobb--Douglas: the flagship example}
For $f(x,y)=x^{\alpha}y^{\beta}$ with $\alpha,\beta>0$ on
$\R^{2}_{++}=\{(x,y):x>0,\ y>0\}$,
\[
  \det\Hess_f \;=\; \alpha\beta\,(1-\alpha-\beta)\,x^{2\alpha-2}y^{2\beta-2},
\]
so (Problem~12)
\[
  f \text{ strictly concave} \iff \alpha+\beta<1, \qquad
  f \text{ concave} \iff \alpha+\beta\le1,
\]
and $f$ is neither convex nor concave when $\alpha+\beta>1$. But $f$ is
\emph{quasiconcave for all} $\alpha,\beta>0$: it is the increasing
transform $\mathrm{e}^{(\cdot)}$ of the strictly concave function
$\alpha\ln x+\beta\ln y$.
\end{rulebox}

\begin{gallery}
\begin{tikzpicture}
  \begin{axis}[surfaxis, title={(a)}, z post scale=1.2]
    \addplot3[surf, domain=-2:2, y domain=-2:2]{x^2+y^2};
  \end{axis}
\end{tikzpicture}\hfill
\begin{tikzpicture}
  \begin{axis}[surfaxis, title={(b)}, z post scale=1.2]
    \addplot3[surf, domain=-2:2, y domain=-2:2]{8-x^2-y^2};
  \end{axis}
\end{tikzpicture}\hfill
\begin{tikzpicture}
  \begin{axis}[surfaxis, title={(c)}, z post scale=1.2]
    \addplot3[surf, domain=-2:2, y domain=-2:2]{x^2-y^2};
  \end{axis}
\end{tikzpicture}\hfill
\begin{tikzpicture}
  \begin{axis}[surfaxis, title={(d)}, z post scale=1.2]
    \addplot3[surf, domain=0.05:2, y domain=0.05:2]{x*y};
  \end{axis}
\end{tikzpicture}
\gcap{(a) $x^{2}+y^{2}$, \;(b) $8-x^{2}-y^{2}$, \;(c) $x^{2}-y^{2}$,
\;(d) $xy$ on the positive quadrant. Problems 1--6 of the additional block
refer to these.}
\end{gallery}

\begin{rulebox}{Why concavity matters: the FOC becomes sufficient}
Let $S\subseteq\R^{n}$ be open and convex and $f$ differentiable on $S$.
\begin{itemize}[leftmargin=1.6em,itemsep=0.3ex]
  \item If $f$ is \emph{concave}, then
        $\grad f(\mathbf{x}^{*})=\mathbf{0}$ $\iff$ $\mathbf{x}^{*}$ is a
        \emph{global} maximizer: stationarity is necessary \emph{and}
        sufficient.
  \item If $f$ is \emph{strictly concave}, it has at most one maximizer,
        so a stationary point is \emph{the unique global maximum}.
  \item Mirror image for convex $f$ and global minima.
  \item Quasiconcavity is too weak for this: $f(x)=x^{3}$ is strictly
        quasiconcave with $f'(0)=0$, yet $0$ is no maximum at all
        (Problem~19). What strict quasiconcavity does guarantee: at most
        one maximizer, and any local maximizer is global.
\end{itemize}
\end{rulebox}

\section{Worked Example}

\begin{worked}{Classify $f(x,y)=x^{2}+xy+y^{2}$ and the Cobb--Douglas $U(x,y)=x^{1/2}y^{1/2}$ on $x,y>0$ --- and see why the second is concave but \emph{not} strictly}
\wstep{Step 1. First function: build the Hessian.}
\[
  f_{x}=2x+y,\quad f_{y}=x+2y,
  \qquad
  \Hess f=\begin{pmatrix}2&1\\1&2\end{pmatrix},
\]
constant, because $f$ is a quadratic.

\wstep{Step 2. Test definiteness by leading principal minors.}
\[
  D_{1}=2>0,\qquad D_{2}=4-1=3>0 .
\]
Positive definite everywhere, so $f$ is \textbf{strictly convex} on all of
$\R^{2}$. Consequently any critical point is the unique global minimum ---
here $(0,0)$.
\wtrap{The matrix is $\begin{pmatrix}2&1\\1&2\end{pmatrix}$, not
$\begin{pmatrix}2&1\\0&2\end{pmatrix}$: the cross term $xy$ contributes
\emph{half} of its coefficient to each off-diagonal entry, because
$f_{xy}=f_{yx}=1$. This is the same halving as in the quadratic forms of
Volume~3, and it changes the answer.}

\wstep{Step 3. Second function: the partials.}
\[
  U_{x}=\tfrac12x^{-1/2}y^{1/2},
  \qquad
  U_{y}=\tfrac12x^{1/2}y^{-1/2} .
\]
Both are positive on $x,y>0$: more of either good is better.

\wstep{Step 4. The second-order partials.}
\[
  U_{xx}=-\tfrac14x^{-3/2}y^{1/2},
  \qquad
  U_{yy}=-\tfrac14x^{1/2}y^{-3/2},
  \qquad
  U_{xy}=\tfrac14x^{-1/2}y^{-1/2} .
\]
$U_{xx}<0$ and $U_{yy}<0$: diminishing marginal utility in each good
separately.

\wstep{Step 5. Now the determinant --- and it vanishes.}
\[
  \det\Hess U=U_{xx}U_{yy}-U_{xy}^{2}
  =\tfrac{1}{16}x^{-1}y^{-1}-\tfrac{1}{16}x^{-1}y^{-1}=0 .
\]
So the Hessian is negative \textbf{semi}definite --- not negative definite.
$U$ is concave, but not \emph{strictly} concave.
\wtrap{This is the chapter's central trap, and the exact-zero result is not
a coincidence: it happens for every Cobb--Douglas whose exponents sum to
$1$. Reporting ``negative definite, so strictly concave'' from
$U_{xx}<0$ alone --- without the determinant --- is the standard error, and
it claims a uniqueness that does not hold.}

\wstep{Step 6. See the flatness directly.}
Along the ray $y=x$, $U(x,x)=x$: a straight line. A strictly concave
function is never linear on any segment, so the failure of strict concavity
is visible without any matrix. It is the constant-returns-to-scale property
of the function --- doubling both inputs doubles output --- expressed as
flatness along rays through the origin.

\wstep{Step 7. What survives, and why economics does not mind.}
$U$ is \textbf{strictly quasiconcave}: its superlevel sets
$\{(x,y):U\ge\bar u\}$ --- the ``at least as good as'' sets --- are strictly
convex, because the indifference curves $xy=\bar u^{2}$ are strictly convex
hyperbolas. Quasiconcavity is weaker than concavity but is what the theory
needs:
\begin{center}\footnotesize
\begin{tabular}{ll}
  \hline
  property & what it buys\\
  \hline
  concave objective on convex set & FOC is sufficient for a global max\\
  \emph{strictly} concave & the maximiser is unique\\
  \emph{strictly quasiconcave} & the \emph{constrained} maximiser is unique\\
  \hline
\end{tabular}
\end{center}
Concavity is not an ordinal property --- $U$ and $U^{3}$ describe the same
preferences while $U^{3}$ is not concave --- so preferences can only ever be
assumed quasiconcave. That is why the textbook assumption is ``convex
preferences'' rather than ``concave utility''.

\wstep{Step 8. If the exponents do not sum to one.}
For $U=x^{a}y^{b}$ with $a,b>0$: the determinant of the Hessian works out
to $ab(1-a-b)x^{2a-2}y^{2b-2}$. So
\[
  a+b<1\Rightarrow\det>0\ \text{and}\ U_{xx}<0:\ \text{strictly concave};
\]
\[
  a+b=1\Rightarrow\det=0:\ \text{concave, not strictly};
  \qquad
  a+b>1\Rightarrow\det<0:\ \textbf{not concave at all.}
\]
Decreasing returns give strict concavity; constant returns give the
borderline case computed above; increasing returns destroy concavity
entirely --- and with it the sufficiency of the first-order conditions,
which is the mathematical reason increasing returns break the competitive
model.
\end{worked}

\section{Practice Problems}

\begin{problems}
  \item Show directly from the chord criterion that an affine function
        $f(x)=ax+b$ is both convex and concave on $\R$, but neither
        strictly.
  \item Show from the chord criterion that $f(x)=|x|$ is convex on $\R$.
        Why is the second-derivative test unavailable here? Is $f$
        strictly convex?
  \item Classify $f(x)=x^{2}-4x+7$ on $\R$ (convex, concave, or neither;
        note strictness).
  \item Classify on $(0,\infty)$: (a) $f(x)=\ln x$; (b) $g(x)=x\ln x$.
  \item Classify $f(x)=x^{3}$, first on $\R$, then on $(0,\infty)$.
  \item For which values of $a\in\R$ is $f(x)=\mathrm{e}^{ax}$ convex on
        $\R$? For which is it strictly convex?
  \item Classify $f(x,y)=x^{2}+y^{2}-4x+6y+13$ on $\R^{2}$ via its
        Hessian, and draw the optimization consequence for its stationary
        point.
  \item Classify $f(x,y)=-x^{2}+xy-y^{2}$ on $\R^{2}$.
  \item Classify $f(x,y)=x^{2}-y^{2}$ on $\R^{2}$.
  \item Classify $f(x,y)=x^{2}+2xy+y^{2}$ on $\R^{2}$. (Careful: this is a
        semidefinite case.)
  \item Classify $f(x,y)=xy$ on $\R^{2}_{++}$.
  \item (Cobb--Douglas, flagship.) Let $f(x,y)=x^{\alpha}y^{\beta}$ with
        $\alpha,\beta>0$ on $\R^{2}_{++}$. Compute the Hessian, show that
        \[
          \det\Hess_f=\alpha\beta(1-\alpha-\beta)\,x^{2\alpha-2}y^{2\beta-2},
        \]
        and classify $f$ according to the value of $\alpha+\beta$.
  \item Using Problem~12, classify on $\R^{2}_{++}$:
        (a) $x^{1/4}y^{1/4}$; \; (b) $x^{1/2}y^{1/2}$; \;
        (c) $x^{2/3}y^{2/3}$.
  \item Show that $f(x,y)=\ln x+\ln y$ is strictly concave on
        $\R^{2}_{++}$, and deduce that $g(x,y)=xy$ is quasiconcave there.
  \item Classify $f(x,y)=-2x^{2}-3y^{2}+4\ln x+6\ln y$ on $\R^{2}_{++}$,
        (a) by the operations rules, (b) directly via the Hessian.
  \item Show that $f(x,y)=\min\{2x+y,\;x+2y\}$ is concave on $\R^{2}$.
        Why does the Hessian test not apply? Is $f$ strictly concave?
  \item Let $g(x,y)=\mathrm{e}^{-x^{2}-y^{2}}$ on $\R^{2}$. Show that $g$
        is \emph{not} concave (evaluate the Hessian at $(1,0)$), but that
        $g$ \emph{is} quasiconcave. Is $g$ quasiconvex?
  \item Show via upper contour sets that $f(x,y)=xy$ is quasiconcave on
        $\R^{2}_{++}$.
  \item Show that $f(x)=x^{3}$ is both quasiconcave and quasiconvex on
        $\R$, and explain what its stationary point $x=0$ teaches about
        FOC sufficiency under mere quasiconcavity.
  \item Classify $f(x,y,z)=x^{2}+2y^{2}+3z^{2}-2xy-2yz$ on $\R^{3}$ using
        the leading principal minors of its Hessian.
  \item For which values of $a\in\R$ is $f(x,y)=ax^{2}+2xy+y^{2}$ convex
        on $\R^{2}$? Strictly convex? Can it ever be concave?
  \item A firm produces output $Q(x,y)=10x+14y-x^{2}-y^{2}-xy$ from inputs
        $x,y$ (prices normalized). Show that $Q$ is strictly concave on
        $\R^{2}$ and find its unique global maximizer and the maximal
        output.
\end{problems}

\section{Solutions}

\begin{solutions}
  \item For any $x,y$ and $\lambda\in[0,1]$:
        \[
          f\bigl(\lambda x+(1-\lambda)y\bigr)
          = a\bigl(\lambda x+(1-\lambda)y\bigr)+b
          = \lambda(ax+b)+(1-\lambda)(ay+b)
          = \lambda f(x)+(1-\lambda)f(y),
        \]
        using $b=\lambda b+(1-\lambda)b$. The chord value \emph{equals}
        the function value, so both defining inequalities hold with
        equality: $f$ is \textbf{convex and concave}. Because equality
        (never strict inequality) holds for all $x\neq y$, it is
        \textbf{neither strictly convex nor strictly concave}.

  \item By the triangle inequality, for $\lambda\in[0,1]$:
        \[
          f\bigl(\lambda x+(1-\lambda)y\bigr)
          = |\lambda x+(1-\lambda)y|
          \;\le\; |\lambda x|+|(1-\lambda)y|
          = \lambda|x|+(1-\lambda)|y| ,
        \]
        since $\lambda,1-\lambda\ge0$. So $f$ is \textbf{convex} on $\R$.
        The second-derivative test is unavailable because $f$ is not
        differentiable at $x=0$ (the graph has a kink). Not strictly
        convex: for $x=1$, $y=2$, $\lambda=\tfrac12$ we get
        $|\tfrac32|=\tfrac32=\tfrac12\cdot1+\tfrac12\cdot2$, equality
        along any chord with both endpoints on the same side of $0$.
        (Remark: being convex, $|x|$ is automatically quasiconvex; it is
        \emph{not} quasiconcave, since for $c>0$ the upper contour set
        $\{|x|\ge c\}=(-\infty,-c]\cup[c,\infty)$ is not convex.)

  \item $f'(x)=2x-4$, $f''(x)=2>0$ for all $x$, so $f$ is
        \textbf{strictly convex} on $\R$. (Consequently its stationary
        point $x=2$ is the unique global minimum, with
        $f(2)=4-8+7=3$.)

  \item (a) $f'(x)=1/x$, $f''(x)=-1/x^{2}<0$ on $(0,\infty)$:
        \textbf{strictly concave}. \;
        (b) $g'(x)=\ln x+1$, $g''(x)=1/x>0$ on $(0,\infty)$:
        \textbf{strictly convex}.

  \item $f''(x)=6x$. On $\R$ this changes sign: $f''(-1)=-6<0$ rules out
        convexity ($f''\ge0$ would be needed everywhere), and $f''(1)=6>0$
        rules out concavity. \textbf{Neither} on $\R$. On $(0,\infty)$:
        $f''(x)=6x>0$, so \textbf{strictly convex} there.

  \item $f''(x)=a^{2}\mathrm{e}^{ax}\ge0$ for every $a$ and every $x$, so
        $f$ is \textbf{convex for all} $a\in\R$. If $a\neq0$ then
        $a^{2}>0$ and $\mathrm{e}^{ax}>0$ give $f''>0$ everywhere:
        \textbf{strictly convex}. If $a=0$ then $f\equiv1$ is affine:
        convex but not strictly.

  \item $f_x=2x-4$, $f_y=2y+6$, so
        $\Hess=\begin{pmatrix}2&0\\0&2\end{pmatrix}$ at every point.
        $D_1=2>0$, $D_2=4>0$: PD everywhere, hence \textbf{strictly
        convex} on $\R^{2}$. Consequence: the stationary point from
        $2x-4=0$, $2y+6=0$, i.e.\ $(2,-3)$, is the \emph{unique global
        minimum}, with $f(2,-3)=4+9-8-18+13=0$.

  \item $f_{xx}=-2$, $f_{yy}=-2$, $f_{xy}=1$:
        $\Hess=\begin{pmatrix}-2&1\\1&-2\end{pmatrix}$ constant.
        $D_1=-2<0$, $D_2=4-1=3>0$: ND everywhere, hence
        \textbf{strictly concave} on $\R^{2}$.

  \item $\Hess=\begin{pmatrix}2&0\\0&-2\end{pmatrix}$ constant.
        $D_2=-4<0$: indefinite at every point, so $f$ is
        \textbf{neither convex nor concave} on $\R^{2}$ (the graph is a
        saddle everywhere).

  \item $f_{xx}=2$, $f_{yy}=2$, $f_{xy}=2$:
        $\Hess=\begin{pmatrix}2&2\\2&2\end{pmatrix}$ constant. Leading
        minors: $D_1=2>0$ but $D_2=4-4=0$, so Sylvester's strict test is
        silent. Check semidefiniteness with \emph{all} principal minors:
        diagonal entries $2\ge0$, $2\ge0$, and $\det=0\ge0$: PSD at every
        point, hence \textbf{convex} on $\R^{2}$. Not strictly: indeed
        $f=(x+y)^{2}$ vanishes identically along the line $y=-x$, so
        chords along that line lie \emph{on} the graph.

  \item $f_{xx}=0$, $f_{yy}=0$, $f_{xy}=1$:
        $\Hess=\begin{pmatrix}0&1\\1&0\end{pmatrix}$ with
        $D_2=-1<0$ at every point of $\R^{2}_{++}$: indefinite everywhere,
        so \textbf{neither convex nor concave}, even restricted to the
        positive orthant. (It is, however, quasiconcave there,
        Problems~14 and~18.)

  \item First partials: $f_x=\alpha x^{\alpha-1}y^{\beta}$,
        $f_y=\beta x^{\alpha}y^{\beta-1}$. Second partials:
        \[
          f_{xx}=\alpha(\alpha-1)x^{\alpha-2}y^{\beta}, \quad
          f_{yy}=\beta(\beta-1)x^{\alpha}y^{\beta-2}, \quad
          f_{xy}=\alpha\beta x^{\alpha-1}y^{\beta-1}=f_{yx}.
        \]
        Determinant:
        \[
          \det\Hess_f
          = \alpha(\alpha-1)\beta(\beta-1)x^{2\alpha-2}y^{2\beta-2}
            -\alpha^{2}\beta^{2}x^{2\alpha-2}y^{2\beta-2}
          = \alpha\beta\,x^{2\alpha-2}y^{2\beta-2}
            \bigl[(\alpha-1)(\beta-1)-\alpha\beta\bigr].
        \]
        Since $(\alpha-1)(\beta-1)-\alpha\beta
        =\alpha\beta-\alpha-\beta+1-\alpha\beta=1-\alpha-\beta$,
        \[
          \det\Hess_f=\alpha\beta(1-\alpha-\beta)\,x^{2\alpha-2}y^{2\beta-2}.
        \]
        The prefactor $\alpha\beta x^{2\alpha-2}y^{2\beta-2}$ is positive
        on $\R^{2}_{++}$, so $\sgn\det\Hess_f=\sgn(1-\alpha-\beta)$.
        Classification:
        \begin{itemize}
          \item $\alpha+\beta<1$: then $\alpha\le\alpha+\beta<1$, so
                $f_{xx}=\alpha(\alpha-1)x^{\alpha-2}y^{\beta}<0$ and
                $\det\Hess_f>0$: ND at every point, hence
                \textbf{strictly concave}.
          \item $\alpha+\beta=1$: then $\alpha<1$ (as $\beta>0$), so
                $f_{xx}<0$, $f_{yy}<0$, and $\det\Hess_f=0$: NSD at every
                point (all principal minors have the right signs), hence
                \textbf{concave, but not strictly}. Indeed $f$ is then
                homogeneous of degree $1$:
                $f(tx,ty)=t^{\alpha+\beta}f(x,y)=tf(x,y)$, so $f$ is
                \emph{linear along every ray} through the origin and
                chords along rays lie on the graph.
          \item $\alpha+\beta>1$: $\det\Hess_f<0$: indefinite everywhere,
                hence \textbf{neither convex nor concave}.
        \end{itemize}

  \item Apply Problem~12. (a) $\alpha+\beta=\tfrac14+\tfrac14=\tfrac12<1$:
        \textbf{strictly concave}. \;
        (b) $\alpha+\beta=\tfrac12+\tfrac12=1$: \textbf{concave, not
        strictly} (linear along rays). \;
        (c) $\alpha+\beta=\tfrac23+\tfrac23=\tfrac43>1$: $\det\Hess<0$
        everywhere, \textbf{neither convex nor concave}. (All three are
        quasiconcave, see the Cobb--Douglas rule and Additional
        Problem~8.)

  \item $f_x=1/x$, $f_y=1/y$, so
        $\Hess=\begin{pmatrix}-1/x^{2}&0\\0&-1/y^{2}\end{pmatrix}$.
        $D_1=-1/x^{2}<0$ and $D_2=1/(x^{2}y^{2})>0$ at every point of
        $\R^{2}_{++}$: ND everywhere, hence \textbf{strictly concave}.
        Now $g(x,y)=xy=\mathrm{e}^{\ln x+\ln y}=F(f(x,y))$ with
        $F=\exp$ increasing. A concave function is quasiconcave, and an
        increasing transform preserves quasiconcavity, so $g$ is
        \textbf{quasiconcave} on $\R^{2}_{++}$: even though Problem~11
        showed it is not concave.

  \item (a) Operations rules: $-2x^{2}$ and $-3y^{2}$ are concave (each is
        a negative multiple of a convex square, and constant in the other
        variable), and $\ln x$, $\ln y$ are (strictly) concave by
        Solution~4(a). A sum with nonnegative weights $2,3,4,6$ of concave
        functions is concave. \;
        (b) Hessian: $f_x=-4x+4/x$, $f_y=-6y+6/y$, so
        \[
          \Hess=\begin{pmatrix}-4-4/x^{2}&0\\0&-6-6/y^{2}\end{pmatrix}.
        \]
        $D_1=-4-4/x^{2}<0$ and
        $D_2=(4+4/x^{2})(6+6/y^{2})>0$ everywhere on $\R^{2}_{++}$: ND,
        hence \textbf{strictly concave}.

  \item Both pieces $2x+y$ and $x+2y$ are affine, hence concave; the
        pointwise minimum of concave functions is concave, so $f$ is
        \textbf{concave} on $\R^{2}$. The two pieces cross where
        $2x+y=x+2y$, i.e.\ along $x=y$; there the gradient jumps from
        $(2,1)$ to $(1,2)$, so $f$ is not differentiable on that line and
        the Hessian test cannot be applied, the min rule delivers the
        verdict anyway. Not strictly concave: on the region $x\le y$
        (where $2x+y\le x+2y$) $f$ equals the affine function $2x+y$, so
        chords there lie on the graph.

  \item Write $g=\mathrm{e}^{h}$ with $h(x,y)=-x^{2}-y^{2}$. Partials:
        $g_x=-2xg$, $g_y=-2yg$, and
        \[
          g_{xx}=(4x^{2}-2)g, \qquad g_{yy}=(4y^{2}-2)g, \qquad
          g_{xy}=4xyg .
        \]
        At $(1,0)$: $g=\mathrm{e}^{-1}$ and
        \[
          \Hess(1,0)=\mathrm{e}^{-1}\begin{pmatrix}2&0\\0&-2\end{pmatrix},
          \qquad D_2=-4\mathrm{e}^{-2}<0 :
        \]
        indefinite, so $g$ is \textbf{not concave} (nor convex) on
        $\R^{2}$. But $h$ has $\Hess_h=-2I$ (ND everywhere), so $h$ is
        strictly concave, hence quasiconcave; and $F=\exp$ is increasing,
        so $g=F\circ h$ is \textbf{quasiconcave}. Check via contour sets:
        for $0<c\le1$,
        $U_c=\{\mathrm{e}^{-x^{2}-y^{2}}\ge c\}
        =\{x^{2}+y^{2}\le-\ln c\}$ is a closed disk (convex); for $c\le0$
        it is all of $\R^{2}$; for $c>1$ it is empty. \emph{Not}
        quasiconvex: for $0<c<1$ the lower contour set
        $L_c=\{x^{2}+y^{2}\ge-\ln c\}$ is the complement of an open disk,
        which is not convex.

  \item Fix $c$. For $c\le0$ the upper contour set is all of
        $\R^{2}_{++}$ (since $xy>0$), which is convex. For $c>0$,
        \[
          U_c=\{(x,y)\in\R^{2}_{++}: y\ge c/x\},
        \]
        the region on or above the graph of $\varphi(x)=c/x$. Since
        $\varphi''(x)=2c/x^{3}>0$, $\varphi$ is convex, and the region
        above the graph of a convex function is a convex set: if
        $y_1\ge\varphi(x_1)$ and $y_2\ge\varphi(x_2)$, then
        \[
          \lambda y_1+(1-\lambda)y_2
          \;\ge\; \lambda\varphi(x_1)+(1-\lambda)\varphi(x_2)
          \;\ge\; \varphi\bigl(\lambda x_1+(1-\lambda)x_2\bigr).
        \]
        Every $U_c$ is convex, so $f$ is \textbf{quasiconcave} on
        $\R^{2}_{++}$, consistent with Solution~14.

  \item $f(x)=x^{3}$ is strictly increasing, so for every $c$ the upper
        contour set is $\{x:x^{3}\ge c\}=[c^{1/3},\infty)$ and the lower
        contour set is $(-\infty,c^{1/3}]$, both intervals, hence
        convex. So $f$ is \textbf{both quasiconcave and quasiconvex}.
        Yet $f'(0)=3\cdot0^{2}=0$: the FOC holds at $x=0$, but $0$ is not
        even a local maximum ($f$ is strictly increasing, and $f$ is
        unbounded above). Moral: for a merely quasiconcave function the
        FOC is \emph{not} sufficient for a maximum, that privilege
        belongs to concavity.

  \item $f_x=2x-2y$, $f_y=4y-2x-2z$, $f_z=6z-2y$, so the (constant)
        Hessian is
        \[
          \Hess=\begin{pmatrix}2&-2&0\\-2&4&-2\\0&-2&6\end{pmatrix}.
        \]
        Leading principal minors: $D_1=2>0$;
        $D_2=2\cdot4-(-2)(-2)=8-4=4>0$;
        \[
          D_3=2\,(4\cdot6-(-2)(-2))-(-2)\,((-2)\cdot6-(-2)\cdot0)+0
             =2\cdot20+2\cdot(-12)=40-24=16>0 .
        \]
        All positive: PD at every point, hence \textbf{strictly convex}
        on $\R^{3}$.

  \item $\Hess=\begin{pmatrix}2a&2\\2&2\end{pmatrix}$ constant. Convexity
        requires PSD, i.e.\ (all principal minors) $2a\ge0$, $2\ge0$,
        $\det=4a-4\ge0$. The binding condition is $a\ge1$ (which implies
        $2a\ge0$): \textbf{convex $\iff a\ge1$}. Strict convexity: PD
        needs $D_1=2a>0$ and $D_2=4a-4>0$, i.e.\ \textbf{$a>1$}. At the
        boundary $a=1$, $f=(x+y)^{2}$ is convex but not strictly
        (Problem~10). For $a<1$, $D_2<0$: indefinite, neither. Concavity
        is impossible for any $a$ because $f_{yy}=2>0$ violates the NSD
        requirement $f_{yy}\le0$.

  \item $Q_x=10-2x-y$, $Q_y=14-2y-x$, so
        $\Hess=\begin{pmatrix}-2&-1\\-1&-2\end{pmatrix}$ constant.
        $D_1=-2<0$, $D_2=4-1=3>0$: ND everywhere, hence $Q$ is
        \textbf{strictly concave} on $\R^{2}$, and any stationary point
        is the unique global maximum. FOC: $10-2x-y=0$ and $14-2y-x=0$.
        From the first, $y=10-2x$; substituting,
        \[
          14-2(10-2x)-x = 14-20+4x-x = 3x-6 = 0
          \;\Rightarrow\; x^{*}=2,\quad y^{*}=10-4=6 .
        \]
        Unique global maximizer $(x^{*},y^{*})=(2,6)$ with
        \[
          Q^{*}=10\cdot2+14\cdot6-4-36-2\cdot6 = 20+84-4-36-12 = 52 .
        \]
\end{solutions}

\section{Additional Problems}

A second pass over the same rules, with fresh functions.

\begin{problems}
  \item Prove directly from the chord criterion that $f(x)=x^{2}$ is
        strictly convex on $\R$: simplify
        $\lambda f(x)+(1-\lambda)f(y)-f\bigl(\lambda x+(1-\lambda)y\bigr)$.
  \item Classify on $(0,\infty)$: (a) $f(x)=\sqrt{x}$; (b) $g(x)=1/x$.
  \item Classify $f(x,y)=3x^{2}-4xy+2y^{2}+x-5y$ on $\R^{2}$.
  \item Classify $f(x,y)=-x^{4}-y^{4}$ on $\R^{2}$. What does this example
        show about the ``ND everywhere'' test for strict concavity?
  \item Classify on $\R^{2}_{++}$: (a) $f(x,y)=x^{0.3}y^{0.5}$;
        (b) $g(x,y)=x^{0.6}y^{0.7}$. Which of the two are quasiconcave?
  \item For which values of $b\in\R$ is $f(x,y)=-x^{2}+bxy-y^{2}$ concave
        on $\R^{2}$? Strictly concave? What happens for $|b|>2$?
  \item Show that the Leontief function $U(x,y)=\min\{x,y\}$ is concave on
        $\R^{2}$, describe its upper contour sets, and confirm
        quasiconcavity.
  \item Show that $f(x,y)=x^{\alpha}y^{\beta}$ is quasiconcave on
        $\R^{2}_{++}$ for \emph{all} $\alpha,\beta>0$.
  \item Classify $f(x,y,z)=-x^{2}-y^{2}-z^{2}+xy$ on $\R^{3}$ using the
        leading principal minors of its Hessian.
  \item A two-product monopolist has profit
        $\pi(q_1,q_2)=20q_1+30q_2-2q_1^{2}-3q_2^{2}-q_1q_2$. Show that
        $\pi$ is strictly concave on $\R^{2}$ and find the unique global
        profit maximizer.

  \item Classify each of the four surfaces (a)--(d) in the figure of the
        Core Rules as convex, concave, both, or neither on the domain
        drawn.
  \item Which of the four are quasiconcave on the domain drawn?
  \item Show by computing the Hessian that $f(x,y)=xy$ is not concave
        anywhere, despite surface (d) looking like a hill from some angles.
  \item Show directly that the superlevel sets of $f(x,y)=xy$ on
        $x,y>0$ are convex, so that $f$ is quasiconcave there.
  \item A function is simultaneously convex and concave on $\R^{n}$. What
        must it be, and what do its level sets look like?
  \item Use the chord criterion to explain in one sentence why a saddle is
        neither convex nor concave.
\end{problems}

\section{Additional Solutions}

\begin{solutions}
  \item Expand the gap between chord and graph:
        \[
          \lambda x^{2}+(1-\lambda)y^{2}
          -\bigl(\lambda x+(1-\lambda)y\bigr)^{2}
          = \lambda x^{2}+(1-\lambda)y^{2}
            -\lambda^{2}x^{2}-2\lambda(1-\lambda)xy-(1-\lambda)^{2}y^{2}.
        \]
        Since $\lambda-\lambda^{2}=\lambda(1-\lambda)$ and
        $(1-\lambda)-(1-\lambda)^{2}=\lambda(1-\lambda)$, this equals
        \[
          \lambda(1-\lambda)x^{2}-2\lambda(1-\lambda)xy
          +\lambda(1-\lambda)y^{2}
          = \lambda(1-\lambda)(x-y)^{2} \;\ge\; 0 ,
        \]
        with strict inequality whenever $x\neq y$ and $\lambda\in(0,1)$.
        Hence $f$ is \textbf{strictly convex} on $\R$.

  \item (a) $f'(x)=\tfrac12x^{-1/2}$,
        $f''(x)=-\tfrac14x^{-3/2}<0$ on $(0,\infty)$:
        \textbf{strictly concave}. \;
        (b) $g'(x)=-x^{-2}$, $g''(x)=2x^{-3}>0$ on $(0,\infty)$:
        \textbf{strictly convex}.

  \item $f_x=6x-4y+1$, $f_y=-4x+4y-5$, so
        $\Hess=\begin{pmatrix}6&-4\\-4&4\end{pmatrix}$ constant.
        $D_1=6>0$, $D_2=24-16=8>0$: PD everywhere, hence
        \textbf{strictly convex} on $\R^{2}$.

  \item $f_x=-4x^{3}$, $f_y=-4y^{3}$, so
        $\Hess=\begin{pmatrix}-12x^{2}&0\\0&-12y^{2}\end{pmatrix}$.
        Principal minors: $-12x^{2}\le0$, $-12y^{2}\le0$,
        $\det=144x^{2}y^{2}\ge0$: NSD at \emph{every} point, hence
        \textbf{concave} on $\R^{2}$. The Hessian is singular wherever
        $x=0$ or $y=0$ (and zero at the origin), so the ``ND
        everywhere'' test cannot certify strictness: yet $f$ \emph{is}
        strictly concave. One variable first: $h(t)=-t^{4}$ has
        $h''(t)=-12t^{2}\le0$ with equality only at the isolated point
        $t=0$, so $h'$ is strictly decreasing and $h$ is strictly
        concave. Then for distinct points
        $(x_1,y_1)\neq(x_2,y_2)$ at least one coordinate differs, and
        \[
          f\bigl(\lambda(x_1,y_1)+(1-\lambda)(x_2,y_2)\bigr)
          = h\bigl(\lambda x_1+(1-\lambda)x_2\bigr)
          + h\bigl(\lambda y_1+(1-\lambda)y_2\bigr)
        \]
        exceeds $\lambda f(x_1,y_1)+(1-\lambda)f(x_2,y_2)$ strictly,
        because the coordinate that differs contributes a strict chord
        inequality and the other contributes at least equality. Lesson:
        ND everywhere is \emph{sufficient, not necessary}, for strict
        concavity.

  \item Apply the Cobb--Douglas rule (Problem~12 of the main list).
        (a) $\alpha+\beta=0.3+0.5=0.8<1$: \textbf{strictly concave}
        (hence also quasiconcave). \;
        (b) $\alpha+\beta=0.6+0.7=1.3>1$:
        $\det\Hess=0.6\cdot0.7\cdot(1-1.3)\,x^{-0.8}y^{-0.6}<0$
        everywhere, so \textbf{neither convex nor concave}: but still
        \textbf{quasiconcave}, by Additional Problem~8. So: both are
        quasiconcave; only (a) is concave.

  \item $\Hess=\begin{pmatrix}-2&b\\b&-2\end{pmatrix}$ constant. NSD
        requires $-2\le0$ (twice, satisfied) and
        $\det=4-b^{2}\ge0$, i.e.\ $|b|\le2$: \textbf{concave $\iff
        |b|\le2$}. Strict concavity: ND needs $D_1=-2<0$ and
        $D_2=4-b^{2}>0$, i.e.\ \textbf{$|b|<2$}. At $b=\pm2$,
        $f=-(x\mp y)^{2}$ is concave but not strictly (it vanishes along
        a whole line). For $|b|>2$, $D_2<0$: indefinite everywhere, so
        $f$ is neither convex nor concave.

  \item The functions $(x,y)\mapsto x$ and $(x,y)\mapsto y$ are affine,
        hence concave; the pointwise minimum of concave functions is
        concave, so $U$ is \textbf{concave} on $\R^{2}$, and therefore
        quasiconcave. Directly: for any $c$,
        \[
          U_c=\{(x,y): \min\{x,y\}\ge c\} = \{x\ge c\}\cap\{y\ge c\}
             = [c,\infty)\times[c,\infty),
        \]
        a ``northeast corner'' set, the intersection of two half-planes,
        hence convex. \; $\checkmark$ (Not strictly concave: on the region
        $x<y$, $U=x$ is affine.)

  \item Take logarithms: $\ln f(x,y)=\alpha\ln x+\beta\ln y$, with
        Hessian
        $\begin{pmatrix}-\alpha/x^{2}&0\\0&-\beta/y^{2}\end{pmatrix}$:
        $D_1=-\alpha/x^{2}<0$, $D_2=\alpha\beta/(x^{2}y^{2})>0$, ND
        everywhere, so $\ln f$ is strictly concave, hence quasiconcave.
        Since $f=\mathrm{e}^{\ln f}$ and $\mathrm{e}^{(\cdot)}$ is
        increasing, and increasing transforms preserve quasiconcavity,
        $f$ is \textbf{quasiconcave} on $\R^{2}_{++}$ for all
        $\alpha,\beta>0$, no restriction on $\alpha+\beta$.

  \item $f_x=-2x+y$, $f_y=-2y+x$, $f_z=-2z$, so the constant Hessian is
        \[
          \Hess=\begin{pmatrix}-2&1&0\\1&-2&0\\0&0&-2\end{pmatrix}.
        \]
        Leading principal minors: $D_1=-2<0$;
        $D_2=(-2)(-2)-1\cdot1=4-1=3>0$; expanding along the third row,
        $D_3=-2\cdot D_2=-6<0$. The pattern $(-,+,-)$ means ND at every
        point, hence \textbf{strictly concave} on $\R^{3}$.

  \item $\pi_{q_1}=20-4q_1-q_2$, $\pi_{q_2}=30-6q_2-q_1$, so
        $\Hess=\begin{pmatrix}-4&-1\\-1&-6\end{pmatrix}$ constant.
        $D_1=-4<0$, $D_2=24-1=23>0$: ND everywhere, hence $\pi$ is
        \textbf{strictly concave} and any stationary point is the unique
        global maximum. FOC: from $20-4q_1-q_2=0$, $q_2=20-4q_1$;
        substituting into $30-6q_2-q_1=0$:
        \[
          30-6(20-4q_1)-q_1 = 30-120+24q_1-q_1 = 23q_1-90 = 0
          \;\Rightarrow\; q_1^{*}=\frac{90}{23},
        \]
        and $q_2^{*}=20-\frac{360}{23}=\frac{460-360}{23}=\frac{100}{23}$.
        Unique global profit maximizer
        $(q_1^{*},q_2^{*})=\bigl(\tfrac{90}{23},\tfrac{100}{23}\bigr)$.

  \item (a) convex --- a bowl. (b) concave --- a dome. (c) neither --- a
        saddle. (d) neither: along the diagonal $y=x$ the function
        $x^{2}$ is convex, while along $y=-x$ it is $-x^{2}$, concave, so
        no single curvature holds.
  \item (a) is quasiconvex but not quasiconcave; (b) is concave, hence
        quasiconcave; (c) is neither; (d) \emph{is} quasiconcave on the
        positive quadrant, which is the whole point of the Cobb--Douglas
        picture --- its superlevel sets are the convex regions beyond the
        isoquants.
  \item $f_{xx}=0$, $f_{yy}=0$, $f_{xy}=1$, so
        $\Hess f=\begin{pmatrix}0&1\\1&0\end{pmatrix}$ with
        $\det=-1<0$: indefinite at every point. An indefinite Hessian rules
        out both concavity and convexity everywhere.
  \item The superlevel set is $\{(x,y):x,y>0,\ xy\ge c\}$. For $c\le0$ it is
        the whole quadrant, which is convex. For $c>0$ it is the region on
        or above the curve $y=c/x$, and that curve is convex as a
        one-variable function ($y''=2c/x^{3}>0$ for $x>0$), so the region
        above it is a convex set. Every superlevel set being convex is the
        definition of quasiconcavity.
  \item It must be \textbf{affine}: $f(\mathbf{x})=\mathbf{a}^{T}
        \mathbf{x}+b$. Both chord inequalities hold with equality, so the
        graph contains every chord, which forces a flat graph --- a
        hyperplane. Its level sets are parallel hyperplanes (parallel
        straight lines when $n=2$).
  \item On a saddle some chords lie above the surface and others lie below
        it --- along one axis the surface curves up, along the other it
        curves down --- so neither the convex nor the concave chord
        inequality can hold for all pairs of points.
\end{solutions}
